\section{本讲小结}
运动描述从参考系、坐标与计时起点开始。位移只由起终点决定，路程还依赖轨迹与往返过程；平均量不能代替瞬时量。速度与加速度是矢量变化率，图像计算应分别识别斜率、有向面积和绝对面积。

当速度随位置给定时，以 $\dd t=\dd x/v$ 累积时间；当位置受几何条件限制时，先写约束再求导。影子、光斑和交点应按几何位置处理。使用转动坐标基矢时，分量和基矢的变化都必须保留。
\section{练习}
\nopagebreak[4]
\begin{question}\begin{drill}[路灯与太阳形成的影子]\label{k1p9}
路灯光源 $S$ 高为 $h$，人高为 $l<h$，人沿水平地面背离灯柱以速度 $v$ 匀速行走。设脚的位置距灯柱为 $x$，头顶影子在地面的位置为 $X$。求影端速度和影长变化率；再讨论太阳光方向固定时的情况。\end{drill}
\fig{k1p9q}{路灯与太阳形成的影子：题设关系}{\begin{tikzpicture}\coordinate(S)at(0,3);\draw[world](0,0)--(6,0);\draw[very thick](0,0)--(S)node[above]{$S$};\draw[world](2.4,0)--(2.4,1.56);\draw[guide](S)--(5,0);\draw[dim](-.5,0)--(-.5,3)node[midway,left]{$h$};\draw[dim](2,0)--(2,1.56)node[midway,left]{$l$};\node[below]at(2.4,0){$x$};\node[below]at(5,0){$X$};\end{tikzpicture}}
\answerlink{k1p9}
\end{question}

\begin{question}\begin{drill}[两圆上部交点的运动]\label{k1p10}
半径均为 $R$ 的两个薄圆环竖直放在同一水平线上。左环圆心 $O_1$ 固定，右环圆心 $O_2$ 沿水平方向运动，$O_1O_2=d$，$0<d<2R$，规定 $\dot d=v$。求两圆上部交点 $P$ 的速度分量和速率。\end{drill}
\fig{k1p10q}{两圆上部交点的运动：题设关系}{\begin{tikzpicture}\draw[world](0,0)circle(1.5) (1.8,0)circle(1.5);\fill(0,0)circle(1pt)node[below]{$O_1$};\fill(1.8,0)circle(1pt)node[below]{$O_2$};\fill(.9,1.2)circle(1pt)node[above]{$P$};\draw[dim](0,-1.9)--(1.8,-1.9)node[midway,below]{$d$};\draw[flow](2,-1.9)--(3,-1.9)node[right]{$v$};\end{tikzpicture}}
\answerlink{k1p10}
\end{question}

\begin{question}\begin{drill}[斜绳牵引的水平速度]\label{k1p11}
物体 $M$ 只能沿水平面向右运动，绳从物体上的 $B$ 点连接到高处固定滑轮 $C$，越过滑轮后以速率 $v$ 收绳。绳不可伸长、滑轮尺寸忽略。当 $BC$ 与水平成 $\alpha$ 时，求物体速度。\end{drill}
\fig{k1p11q}{斜绳牵引的水平速度：题设关系}{\begin{tikzpicture}\draw[world](0,0)--(5,0);\draw[world](.2,0)rectangle(1,.5);\node at(.6,.25){$M$};\draw[world](1,.5)--(4,2.5)--(5,2.5);\draw[world](4,2.5)circle(.12);\node[above]at(4,2.6){$C$};\node[below right]at(1,.5){$B$};\node at(1.5,.65){$\alpha$};\draw[flow](4.5,2.9)--(5.3,2.9)node[right]{$v$};\end{tikzpicture}}
\answerlink{k1p11}
\end{question}

\begin{question}\begin{drill}[极坐标速度的矢量求导]\label{k1p12}
在固定平面内，质点的极坐标速度写成 $\bm v=v_\rho\bm e_\rho+\rho\omega\bm e_\theta$，其中 $v_\rho=\dot\rho$、$\omega=\dot\theta$。推导加速度在径向与横向的分量。\end{drill}
\fig{k1p12q}{极坐标速度的矢量求导：题设关系}{\begin{tikzpicture}\draw[axis](0,0)--(4,0)node[right]{$x$};\draw[axis](0,0)--(0,3)node[above]{$y$};\draw[guide](0,0)--(2.5,1.5)node[midway,above]{$\rho$};\fill(2.5,1.5)circle(1pt)node[below]{$P$};\draw[flow](2.5,1.5)--(3.6,2.16)node[right]{$\bm e_\rho$};\draw[flow](2.5,1.5)--(1.9,2.5)node[above]{$\bm e_\theta$};\node at(.8,.25){$\theta$};\end{tikzpicture}}
\answerlink{k1p12}
\end{question}
